\documentclass{article}

% Language setting
% Replace `english' with e.g. `spanish' to change the document language
\usepackage[english]{babel}

% Set page size and margins
% Replace `letterpaper' with `a4paper' for UK/EU standard size
\usepackage[letterpaper,top=2cm,bottom=2cm,left=3cm,right=3cm,marginparwidth=1.75cm]{geometry}

% Useful packages
\usepackage{amsmath}
\usepackage{graphicx}
\usepackage[colorlinks=true, allcolors=blue]{hyperref}

\title{Time Series Coursework}
\author{Gianluca Fiorini, Tommaso Di Bene, Riccardo Visentini, Nikola Stevanovic}

\begin{document}
\maketitle

\section{Introduction \& Objective}
Choose two time series from the fMRI dataset that we have discussed in class and that you will find in Virtuale and further details on the dataset are also attached. Choose the two time series in such a way that one of them is non Gaussian and the other is (you can test for normality or you can detect non-Gaussianity in the diagnostic check).
Fit an appropriate score-driven model and compare it with a Gaussian model. Discuss the outputs of the analysis.

\section{Dataset Description}
The fMRI dataset contains information about dynamic functional activity, monitored for each brain region in the Desikan atlas, which is gathered from resting state fMRI (R-fMRI).\\\\
This imaging technology monitors brain functional activity at different regions via dynamic changes in blood flow creating a low frequency blood-oxygen-level dependent (BOLD) signal – when the subject is not performing an explicit task during the imaging session.\\\\
In the NKI1 study, the subjects were asked to stay awake with eyes open.
As for the structural networks, the raw R-fMRI scans are pre-processed to derive the time-series data for each brain region. In this case the C-PAC software has been used. C-PAC lets you easily export BOLD timeseries in a number of different ways. This can be useful for those wishing to undertake advanced analysis not included in C-PAC, but still take advantage of its robust pre-processing features.\\\\
The dataset on dynamic functional activity contains an array Y of dimensions 70×404×24×2 comprising the 70×404 multivariate timeseries activity data collected for the 24 subjects in each of the 2 scan-rescan imaging sessions. In particular Y[,,i,k] is a 70 × 404 matrix whose rows contain the dynamic activity data of the brain regions, collected at T = 404 equally spaced times, with time lag between consecutive time observations being equal to 1400 ms, for subject i, monitored during scan k, for every i = 1, . . . , 24 and k = 1, 2.\\\\
Based on this description, the element Y[v,t,i,k] ∈ R is a real number
measuring the activity of brain region v at time t, for subject i, during scan k.\\\\
Dynamic functional activity data for 2 subjects are not available, whereas for other 10 subjects only the first scan is observed. Therefore, their corresponding entries in Y are empty.\\\\
Finally, during pre-processing, the spatial dependence has not been filtered out, and hence may play an important role in joint modeling of such data.\\\\
From the time-series data, characterizing dynamic functional activity,
the functional networks, measuring synchronization in activity for each pair of brain regions, are also derived by simply computing the Pearson correlation coefficient for each pair of time-series data. These information are contained in the 70×70×24×2 array W which can be found in the file fMRI-connectome.RData. In particular W[,,i,k] is a 70×70 symmetric matrix measuring pairwise correlations for each pair of brain regions in subject i, monitored during scan k, for every i = 1, . . . , 24 and k = 1, 2.\\\\
Based on this description, the element W[v,u,i,k] ∈ [−1, 1] is simply equal to cor(Y[v,,i,k],Y[u,,i,k]). As can be noticed by a graphical analysis of these correlation matrices, some of them have an overall much higher functional correlation than others. This behavior has been noticed in the past by the computational neuroscientists who provided the data, with different pipelines, and their suggestion is to convert correlations to ranks. However, developing more consistent measures of functional synchronization or statistical models robust to this scaling effect may be relevant research directions.


\subsection{This could be included in a methodology section}
The Configurable Pipeline for the Analysis of Connectomes (C-PAC) is an open-source software pipeline for automated preprocessing and analysis of resting-state fMRI data. C-PAC builds upon a robust set of existing software packages including AFNI, FSL, and ANTS, and makes it easy for both novice users and experts to explore their data using a wide array of analytic tools. Users define analysis pipelines by specifying a combination of preprocessing options and analyses to be run on an arbitrary number of subjects. Results can then be compared across groups using the integrated group statistics feature.

Optimized For Large Datasets
Researchers can now explore functional connectome data from thousands of subjects made public through releases by the International Neuroimaging Data-Sharing Initiative (INDI) and 1000 Functional Connectomes Project. With this in mind, C-PAC has been designed to reliably preprocess and analyize data for hundreds of subjects in a single run, either on a single machine or on a compute cluster using Sun Grid Engine, HTCondor, or Portable Batch System.

Ensure Robust Results
Different analysis pipelines can produce significantly different results, raising questions about the replicability and reliability of brain imaging findings. C-PAC makes it easy to explore the impact of particular processing decisions by allowing users to run a factorial number of analysis pipelines, each with a different set of preprocessing and analysis options. An integrated Quality Control interface facillitates rapid manual examination of pipeline outputs and selection of subjects to be included in group comparisons.\\\\




\subsection{How to include Figures}

First you have to upload the image file from your computer using the upload link in the file-tree menu. Then use the includegraphics command to include it in your document. Use the figure environment and the caption command to add a number and a caption to your figure. See the code for Figure \ref{fig:frog} in this section for an example.

Note that your figure will automatically be placed in the most appropriate place for it, given the surrounding text and taking into account other figures or tables that may be close by. You can find out more about adding images to your documents in this help article on \href{https://www.overleaf.com/learn/how-to/Including_images_on_Overleaf}{including images on Overleaf}.

\begin{figure}
\centering
\includegraphics[width=0.25\linewidth]{frog.jpg}
\caption{\label{fig:frog}This frog was uploaded via the file-tree menu.}
\end{figure}

\subsection{How to add Lists}

You can make lists with automatic numbering \dots

\begin{enumerate}
\item Like this,
\item and like this.
\end{enumerate}
\dots or bullet points \dots
\begin{itemize}
\item Like this,
\item and like this.
\end{itemize}

\subsection{How to write Mathematics}

\LaTeX{} is great at typesetting mathematics. Let $X_1, X_2, \ldots, X_n$ be a sequence of independent and identically distributed random variables with $\text{E}[X_i] = \mu$ and $\text{Var}[X_i] = \sigma^2 < \infty$, and let
\[S_n = \frac{X_1 + X_2 + \cdots + X_n}{n}
      = \frac{1}{n}\sum_{i}^{n} X_i\]
denote their mean. Then as $n$ approaches infinity, the random variables $\sqrt{n}(S_n - \mu)$ converge in distribution to a normal $\mathcal{N}(0, \sigma^2)$.


\subsection{How to change the margins and paper size}

Usually the template you're using will have the page margins and paper size set correctly for that use-case. For example, if you're using a journal article template provided by the journal publisher, that template will be formatted according to their requirements. In these cases, it's best not to alter the margins directly.

If however you're using a more general template, such as this one, and would like to alter the margins, a common way to do so is via the geometry package. You can find the geometry package loaded in the preamble at the top of this example file, and if you'd like to learn more about how to adjust the settings, please visit this help article on \href{https://www.overleaf.com/learn/latex/page_size_and_margins}{page size and margins}.

\subsection{How to change the document language and spell check settings}

Overleaf supports many different languages, including multiple different languages within one document.

To configure the document language, simply edit the option provided to the babel package in the preamble at the top of this example project. To learn more about the different options, please visit this help article on \href{https://www.overleaf.com/learn/latex/International_language_support}{international language support}.

To change the spell check language, simply open the Overleaf menu at the top left of the editor window, scroll down to the spell check setting, and adjust accordingly.

\subsection{How to add Citations and a References List}

You can simply upload a \verb|.bib| file containing your BibTeX entries, created with a tool such as JabRef. You can then cite entries from it, like this: \cite{greenwade93}. Just remember to specify a bibliography style, as well as the filename of the \verb|.bib|. You can find a \href{https://www.overleaf.com/help/97-how-to-include-a-bibliography-using-bibtex}{video tutorial here} to learn more about BibTeX.

If you have an \href{https://www.overleaf.com/user/subscription/plans}{upgraded account}, you can also import your Mendeley or Zotero library directly as a \verb|.bib| file, via the upload menu in the file-tree.

\subsection{Good luck!}

We hope you find Overleaf useful, and do take a look at our \href{https://www.overleaf.com/learn}{help library} for more tutorials and user guides! Please also let us know if you have any feedback using the \textbf{Contact us} link at the bottom of the Overleaf menu --- or use the contact form at \url{https://www.overleaf.com/contact}.

\bibliographystyle{alpha}
\bibliography{sample}

\end{document}
